% Einheit 3 – Kreisteile – Nr. 33–43
\setcounter{aufgabe}{32}
\einheitenkopf[e3][3 Kreisteile]{Einheit 3 von 3 · Kreisteile}
\zweigzeile{Hier lernst du, Anteil, Bogen und Fläche eines Kreisausschnitts aus dem Mittelpunktswinkel zu berechnen · neu oder schon bekannt – je nach Buch, manche Bücher erst Kl. 9 · P10 · baut auf: Kreisumfang (Einheit 1), Kreisfläche (Einheit 2), Winkel, Anteile als Bruch und Prozent}

\begin{aufgabe}{Ich erkenne, welcher Teil vom Kreis gemeint ist.}
Schreibe auf, welcher Teil vom Kreis grau ist: Halbkreis, Viertelkreis, Dreiviertelkreis, Achtelkreis oder Drittelkreis.

\begin{kreis}[1]\mittelpunkt{}\sektor{0}{180}{}\end{kreis}\hfill
\begin{kreis}[1]\mittelpunkt{}\sektor{90}{180}{}\end{kreis}\hfill
\begin{kreis}[1]\mittelpunkt{}\sektor{0}{270}{}\end{kreis}\hfill
\begin{kreis}[1]\mittelpunkt{}\sektor{0}{45}{}\end{kreis}\hfill
\begin{kreis}[1]\mittelpunkt{}\sektor{90}{210}{}\end{kreis}

\begin{teilezwei}
\tz 1. Figur: \leerfeld & \tz 2. Figur: \leerfeld \\
\tz 3. Figur: \leerfeld & \tz 4. Figur: \leerfeld \\
\tz 5. Figur: \leerfeld & \\
\end{teilezwei}
\end{aufgabe}

\verfahren{Anteil und Mittelpunktswinkel}
\begin{aufgabe}{Ich kann den Anteil eines Kreisausschnitts am ganzen Kreis angeben.}
Der Mittelpunktswinkel $\alpha$ ist gegeben. Gib den Anteil am ganzen Kreis als gekürzten Bruch an.
\begin{teilezwei}
\tz $\alpha = 90^\circ$: \leerfeld & \tz $\alpha = 180^\circ$: \leerfeld \\
\tz $\alpha = 60^\circ$: \leerfeld & \tz $\alpha = 120^\circ$: \leerfeld \\
\anweisung{Gib den Anteil in Prozent an. Runde auf eine Stelle nach dem Komma.}
\tz $\alpha = 36^\circ$: \leerfeld[\%] & \tz $\alpha = 100^\circ$: \leerfeld[\%] \\
\end{teilezwei}
\begin{teile}
\teil Auf einem Glücksrad ist der graue Bereich der Gewinn. Der weiße Bereich hat den Mittelpunktswinkel $125^\circ$. Wie viel Prozent des Glücksrads sind grau? Runde auf eine Stelle. (P10 2025 OS) \hfill \leerfeld[\%]
\end{teile}

\begin{kreis}[1.3]\mittelpunkt{}\sektor{0}{235}{}\mittelpunktswinkel{235}{360}{125^\circ}\end{kreis}

\end{aufgabe}

\begin{aufgabe}{Ich kann aus dem Anteil den Mittelpunktswinkel berechnen.}
Berechne den Mittelpunktswinkel $\alpha$ des Ausschnitts.
\begin{teilezwei}
\tz Anteil $\frac{1}{4}$: \quad \leerfeld[°] & \tz Anteil $\frac{2}{3}$: \quad \leerfeld[°] \\
\tz Anteil $30\,\%$: \quad \leerfeld[°] & \\
\end{teilezwei}
\end{aufgabe}

\verfahren{Bogen, Ausschnitt und Ring}
\begin{aufgabe}{Ich kann die Länge eines Kreisbogens berechnen.}
Berechne die Länge des Kreisbogens zum Mittelpunktswinkel $\alpha$. Runde auf eine Stelle nach dem Komma.
\begin{teile}
\teil $\alpha = 90^\circ$, \ $r = 4\,\text{cm}$ \quad \leerfeld[cm]
\teil $\alpha = 60^\circ$, \ $r = 9\,\text{cm}$ \quad \leerfeld[cm]
\teil $\alpha = 135^\circ$, \ $r = 8\,\text{cm}$ \quad \leerfeld[cm]
\teil $\alpha = 250^\circ$, \ $d = 5\,\text{cm}$ \quad \leerfeld[cm]
\end{teile}
\end{aufgabe}

\begin{aufgabe}{Ich kann die Fläche eines Kreisausschnitts berechnen.}
Berechne die Fläche des Kreisausschnitts. Runde auf eine Stelle nach dem Komma.
\begin{teile}
\teil $\alpha = 90^\circ$, \ $r = 4\,\text{cm}$ \quad \leerfeld[cm²]
\teil $\alpha = 120^\circ$, \ $r = 6\,\text{cm}$ \quad \leerfeld[cm²]
\teil $\alpha = 200^\circ$, \ $r = 7{,}5\,\text{cm}$ \quad \leerfeld[cm²]
\teil $\alpha = 300^\circ$, \ $d = 9\,\text{m}$ \quad \leerfeld[m²]
\end{teile}
\end{aufgabe}

\begin{aufgabe}{Ich kann den Umfang eines Kreisausschnitts berechnen.}
Berechne die Länge des ganzen Randes: Bogen und die beiden Radien. Runde auf eine Stelle nach dem Komma.

\begin{kreis}[1.2]\mittelpunkt{M}\sektor{0}{100}{\alpha}\end{kreis}

\begin{teile}
\teil $\alpha = 90^\circ$, \ $r = 4\,\text{cm}$ \quad \leerfeld[cm]
\teil $\alpha = 60^\circ$, \ $r = 12\,\text{cm}$ \quad \leerfeld[cm]
\teil $\alpha = 150^\circ$, \ $r = 5\,\text{cm}$ \quad \leerfeld[cm]
\end{teile}
\end{aufgabe}

\begin{aufgabe}{Ich kann die Fläche eines Kreisrings berechnen.}
Zwei Kreise haben denselben Mittelpunkt. Berechne die Fläche des Rings zwischen ihnen. Runde auf eine Stelle nach dem Komma.
\begin{teile}
\teil äußerer Radius $6\,\text{cm}$, innerer Radius $4\,\text{cm}$ \quad \leerfeld[cm²]
\teil äußerer Durchmesser $30\,\text{cm}$, innerer Durchmesser $20\,\text{cm}$ \quad \leerfeld[cm²]
\teil äußerer Radius $9\,\text{cm}$, der Ring ist $3\,\text{cm}$ breit \quad \leerfeld[cm²]
\end{teile}
\end{aufgabe}

\begin{aufgabe}{Ich finde den Fehler bei Kreisteilen.}
Finde den Fehler, benenne ihn und korrigiere die Rechnung.
\begin{teile}
\teil Anteil des Ausschnitts mit $\alpha = 40^\circ$: \rechnung{40 : 180 &= \tfrac{2}{9}}
\teil Umfang des Ausschnitts mit $\alpha = 90^\circ$ und $r = 6\,\text{cm}$: \rechnung{u &= \tfrac{1}{4} \cdot 2 \cdot \pi \cdot 6\,\text{cm} \\ &\approx 9{,}4\,\text{cm}}
\teil Der weiße Teil hat $110^\circ$. Wie viel Prozent des Kreises sind grau? \rechnung{110 : 360 &\approx 0{,}306 = 30{,}6\,\%}
\end{teile}

\begin{kreis}[1.1]\mittelpunkt{}\sektor{0}{250}{}\mittelpunktswinkel{250}{360}{110^\circ}\end{kreis}

\end{aufgabe}

\begin{aufgabe}{Ich kann begründen, wie Mittelpunktswinkel und Anteil zusammenhängen.}
Begründe.
\begin{teile}
\teil Warum gehört ein Halbkreis zum Mittelpunktswinkel $180^\circ$?
\teil Zwei Kreisausschnitte haben beide den Mittelpunktswinkel $60^\circ$, aber verschiedene Radien. Haben sie denselben Anteil an ihrem Kreis? Haben sie dieselbe Fläche?
\end{teile}
\end{aufgabe}

\begin{aufgabe}{Ich kann mit einem Kreisausschnitt eine Rasenfläche berechnen.}
Ein Rasensprenger spritzt $7\,\text{m}$ weit und dreht sich dabei um $150^\circ$ hin und her.

\begin{kreis}[1.6]\mittelpunkt{}\sektor{0}{150}{150^\circ}\radius{0}{7\,\text{m}}\end{kreis}

\begin{teile}
\teil Wie groß ist die Fläche, die er bewässert? Runde auf eine Stelle. \hfill \leerfeld[m²]
\teil Wie lang ist der äußere, gebogene Rand dieser Fläche? Runde auf eine Stelle. \hfill \leerfeld[m]
\teil Der Sprenger soll $80\,\text{m}^2$ bewässern. Um wie viel Grad muss er sich drehen? Runde auf ganze Grad. \hfill \leerfeld[°]
\end{teile}
\end{aufgabe}

\begin{aufgabe}{Ich kann ein Kreisdiagramm berechnen und zeichnen.}
Von 40 Kindern einer Klasse kommen 18 mit dem Rad, 12 mit dem Bus und 10 zu Fuß.
\begin{teile}
\teil Berechne die Anteile in Prozent. \\ Rad \leerfeld[\%] \quad Bus \leerfeld[\%] \quad zu~Fuß \leerfeld[\%]
\teil Berechne die Mittelpunktswinkel. \\ Rad \leerfeld[°] \quad Bus \leerfeld[°] \quad zu~Fuß \leerfeld[°]
\teil Zeichne das Kreisdiagramm in den Kreis und beschrifte die Teile.
\end{teile}

\kreisleer[2.2]

\end{aufgabe}
